\section{A direct comparison for small noise}
\label{v10:smallnoise}

The small-noise part of the comparison can be proved from the incoming
$0.6$ strip itself. Neither the sharper half-amplitude bound nor the
fractional-part and marked-crossing estimates are needed in the following
range. This is a comparison of actual conditional sum laws; contact and a
predecessor bound are needed only to obtain its premises elsewhere.

\begin{theorem}\label{v10:smallnoise:theorem}
Let $n\ge42$, $2/5\le p\le53/100$, $q=1-p$, and $Y=B+U$, where
$\Pp(B=1)=p$ and $\E(U\mid B)=0$. Suppose the two conditional laws satisfy
\[
 |U_j|\le3/5,\qquad U_0\le a_0,\qquad U_1\ge-a_1,
 \qquad a_0,a_1\ge0,\quad a_0+a_1\le9/125.
\]
Put $v=pq$, $s=\E U^2$, $\lambda=ns$, $\ell=(v+s)^{-1/2}$,
$X=\ell(Y-p)$, and $\rho=\E|X|^3$. Assume
$0<\lambda\le1/500$ and $|H|\le46/25$.
Choose an integer $k$ and $u\in[-34/125,91/125]$ such that
$np+H/\ell=k+u$. Write $\delta=k-np$, $\rho_B=(1-2v)/\sqrt v$,
\[
 a=\sqrt n\{\Pp(S_n\le H)-\Phi(H/\sqrt n)\}-\CE\rho,
 \quad
 C_B=\frac{\sqrt n}{\rho_B}
       \{F_B(k)-\Phi(\delta/\sqrt{nv})\}.
\]
Then $1\le k\le n-1$ and
\[
                 a+\rho_B(\CE-C_B)<-\lambda/100.
\]
In particular no law in this range can have $a\ge0$, by the sharp
Bernoulli bound. The theorem does not cover larger noise or the remaining
probability panels.
\end{theorem}

\subsection{The exact comparison and uniform normalization}
Let $W_j$ have the law of the sum of $j$ independent $U_1$ variables and
$n-j$ independent $U_0$ variables, and put
$\pi_j=\binom njp^jq^{n-j}$, zero outside the legal counts. Define
\[
 f(u)=\frac{\Phi((\delta+u)/\sqrt{n(v+s)})-
                   \Phi(\delta/\sqrt{n(v+s)})}{\pi_k},
 \quad O=\Pp(W_k>u)+f(u),
 \quad T=\sum_{i=1}^{n-k}\Pp(W_{k+i}\le u-i).
\]
Conditioning on the sign count, including its atoms, gives
\[
 Q:=\Pp(S_n\le H)-F_B(k)
 =-\sum_{i=0}^k\pi_{k-i}\Pp(W_{k-i}>u+i)
   +\sum_{i=1}^{n-k}\pi_{k+i}\Pp(W_{k+i}\le u-i).
\]
Consequently, if $\pi_{k+i}/\pi_k\le R_0$,
\begin{equation}\label{v10:smallnoise:comparison}
 a+\rho_B(\CE-C_B)
 \le-\sqrt n\pi_k(O-R_0T)+\mathcal R,
 \quad
 \mathcal R=\sqrt n\{\Phi(\delta/\sqrt{nv})-
                  \Phi(\delta/\sqrt{n(v+s)})\}+\CE(\rho_B-\rho).
\end{equation}
The discarded terms are nonpositive. The outgoing event is strict and
the incoming events closed, also at a collision.

Here $v\ge6/25$ and
$\sqrt{v+s}<501/1000$. Thus $|\delta|<33/20$ throughout;
on $|u|\le1/20$ it is less than one. These inequalities imply the stated
interior-count range and $9/25<k/n<23/40$.
On this last interval the absolute third derivative of binary relative
entropy is at most $6$, the absolute logarithmic derivative of $r(1-r)$
at most $5/4$, and $1/(12r(1-r))<3/8$. The finite Stirling estimate
therefore gives
\[
 \sqrt n\pi_k=\frac{\phi(0)}{\sqrt v}
       \exp\{-\delta^2/(2nv)+e\},\qquad
 n|e|<Z(D):=5D/8+3/8+D^3/42
\]
whenever $|\delta|\le D\le33/20$.
This follows by expanding relative entropy to second order and the
logarithmic prefactor to first order; the three remainders are respectively
$D^3/n^2$, $5D/(8n)$ and $3/(8n)$. It uses the same finite Stirling
formula proved in the manuscript's count appendix, not an asymptotic law.
Using $0.39886<\phi(0)<0.4$, $e^{-x}\ge1-x$, and
$e^x\le(1-x)^{-1}$ for $x<1$ yields the following bounds:
\[
\begin{array}{c|cc}
 & |u|\le1/20 & -34/125\le u\le91/125\\\hline
 \sqrt n\pi_k\text{ lower bound}&73/100&13/20\\
 f'\text{ lower bound}&19/20&4/5\\
 f'\text{ upper bound}&21/20&5/4\\
 \pi_{k+i}/\pi_k\text{ upper bound}&53/50&3/2\\
 \mathcal R\text{ upper bound}&2\lambda/n&3\lambda/n
\end{array}
\]
In the derivative bounds the actual $\delta$ is held fixed and the
derivative variable runs between zero and the actual phase. Indeed
\[
 \log f'(t)=-\tfrac12\log(1+s/v)+\frac{\delta^2}{2nv}
                  -\frac{(\delta+t)^2}{2n(v+s)}-e,
\]
and, for $|t|\le b$,
\[
 n|\log f'(t)|\le(8/3)\lambda+
       \frac{2Db+b^2}{2(6/25)}+Z(D).
\]
Take $(D,b)=(1,1/20)$ or $(33/20,91/125)$.
At a binomial mode, $|\delta|\le1$; the same Stirling estimate shows
$\sqrt n\max_j\pi_j<17/20$. Dividing this by $13/20$ proves the
global ratio bound. In the inner phase range, the exact binomial recursion
gives $\pi_{k+1}/\pi_k\le1+(53/100)/(42\cdot6/25)<53/50$;
all subsequent factors are at most one since $(n+1)p-k<153/100<2$.

For the repair bound, conditional signs are fixed because $a_0<p$ and
$a_1<q$. Expansion of the absolute cubic gives, with $r=\min(p,q)$,
\[
 \E|Y-p|^3\ge v(1-2v)+(3r-d)s,\qquad d=a_0+a_1.
\]
This follows term by term from
$U_1^3\ge-a_1U_1^2$ and $U_0^3\le a_0U_0^2$.
Since $(1+x)^{3/2}\le1+3x/2+3x^2/8$, it follows that
\[
 (v+s)^{3/2}(\rho-\rho_B)
 \ge s\{3r-\tfrac32(1-2v)-d-3(1-2v)s/(8v)\}>0.
\]
The bracket is positive using $r\ge2/5$, $v\ge6/25$ and
$s\le1/21000$. We may therefore discard the favorable moment term.
Differentiating the normal term in $s$ bounds the remaining repair by
$(7/4)\delta_+s$, because $0.2/(6/25)^{3/2}<7/4$. This proves the two
repair bounds in the table.

\subsection{Incoming counts need only a fourth moment}
For a centered independent sum $W$ whose summands have absolute bound
$A=3/5$ and total variance $t$,
\[
                  \E W^4\le A^2t+3t^2.
\]
Also $|s_1-s_0|\le(5/2)s$, so that
$t_{k+i}\le iV\lambda$, with $V=9/8$ in the inner phase range and
$V=7/6$ globally. Markov's inequality, applied to a closed incoming event,
therefore gives, for any phase upper bound $h\in[0,1)$,
\begin{equation}\label{v10:smallnoise:tails}
 \frac{T}{\lambda}\le\frac{K(V)}{(1-h)^4},\qquad
 K(V)=\frac9{25}V\frac{49}{40}
                         +3V^2\frac1{500}\frac{17}{10}.
\end{equation}
Indeed $i-u\ge i(1-h)$; sum the resulting $i^{-3}$ and $i^{-2}$ terms.
The elementary integral bounds
\[
 \sum_{i\ge1}i^{-3}\le1+1/8+1/27+1/18<49/40,
 \quad
 \sum_{i\ge1}i^{-2}\le1+1/4+1/9+1/3<17/10
\]
justify these constants. For $u\le0$ use $h=0$.

\subsection{The own count, uniformly down to zero noise}
Set $\alpha=9/20$ and
\[
 h_\alpha(x)=\frac{-x^4+3\alpha^2x^2+2\alpha^3x}{5\alpha^4}.
\]
Then $h_\alpha\le\mathbf1_{(0,\infty)}$: the polynomial is nonpositive
on the negative half-line and its positive maximum is
$(9+6\sqrt3)/20<1$. Its value at zero is zero. For a centered independent
sum $W$ as above, Cauchy and the moment expansion give
$|\E W^3|\le At$ and $\E W^4\le A^2t+3t^2$. Thus
\[
 \Pp(W+\sigma r>0)\ge
 \frac{(3\alpha^2-A^2)t-3t^2-4Art-6r^2t-r^4
                         +3\alpha^2r^2+2\sigma\alpha^3r}{5\alpha^4}.
\]
In the inner phase range,
$79\lambda/84\le t=\E W_k^2\le89\lambda/84$.
Put $T_*=89/42000$. To bound the preceding expression plus $br$, divide
by $t>0$ and write $r=zt$. The result is
\[
 F_{\sigma,b}(z,t)=\frac{3\alpha^2-A^2-3t}{5\alpha^4}
 +\left(\frac{2\sigma}{5\alpha}+b-\frac{4At}{5\alpha^4}\right)z
 +\frac{3\alpha^2t-6t^2}{5\alpha^4}z^2
 -\frac{t^3z^4}{5\alpha^4}.
\]
It is concave in $t\ge0$. On $0\le z\le z_0$, replace $z^4$ by
$z_0^2z^2$ to obtain a lower quadratic. The minimum of this quadratic
on the interval occurs at an endpoint or its eligible vertex. Checking
it at $t=0,T_*$ proves
\[
 F_{1,-21/20}(z,t)>4/5\quad(0\le z\le2),\qquad
 F_{-1,19/20}(z,t)>1\quad(0\le z\le3/2).
\]
The respective rational margins are greater than $0.0288$ and $0.176$.
These are interval inequalities, not sampled noise checks.

For negative phase $u=-r$ and $r\ge2t$, Cantelli instead gives
$O\ge r^2/(t+r^2)-(21/20)r$. At $r=2t$ its value divided by $t$
exceeds $4/5$, and its derivative in $r$ is positive. That derivative
first increases and then decreases, so any subsequent stationary point
is a maximum. The value at the other endpoint $r=1/20$ also exceeds
$(4/5)t$, checked at $t=T_*$. Hence $O>(4/5)t$ throughout the negative
inner phase range. For positive phase, the quartic handles $r\le3t/2$;
for larger $r$, simply use $O\ge f(r)\ge(19/20)r>t$.
Both switching points are below $1/20$. We have proved
\[
                         O>\frac45\frac{79}{84}\lambda
                         \qquad (|u|\le1/20).
\]
Combining this with \eqref{v10:smallnoise:tails} and the table yields
\[
 \frac{73}{100}\left\{\frac45\frac{79}{84}
       -\frac{53}{50}\frac{K(9/8)}{(19/20)^4}\right\}
                      -\frac2{42}-\frac1{100}
   =\frac{219609181}{27367410000}>0.
\]
The bracket is positive before the lower mass bound is used in
\eqref{v10:smallnoise:comparison}.

\subsection{The two outer phase intervals}
For $u=-r\in[-34/125,-1/20]$, Cantelli and the global derivative bound give
\[
 O\ge\frac{r^2}{11/5000+r^2}-\frac54r.
\]
This expression is concave on the stated interval because
$3r^2>11/5000$. Its smaller endpoint value exceeds $0.469$.
Subtract $(3/2)K(7/6)/500$ for incoming counts. Multiplication by the
mass lower bound $13/20$ still exceeds the repair $3/(42\cdot500)$
and the desired $1/50000$ margin.

For $u\in[1/20,91/125]$ use only $O\ge(4/5)u$.
The function $u(1-u)^4$ increases and then decreases on $[0,1]$;
its minimum on this interval is
$G=(91/125)(34/125)^4$. Consequently
\[
 O-R_0T\ge (4/5)G-(3/2)K(7/6)/500>0.
\]
Here the positive numerator was divided by $(1-u)^4\le1$ before
dropping that denominator. Finally
\[
 \frac{13}{20}\{(4/5)G-(3/2)K(7/6)/500\}
       -\frac3{42\cdot500}-\frac1{50000}
       =\frac{4806449287359}{5468750000000000}>0.
\]
Using the upper noise endpoint is valid in both outer intervals since
every discarded error increases with $\lambda$.
This completes the proof of Theorem~\ref{v10:smallnoise:theorem}.

The accompanying exact rational checker reconstructs all 46 scalar gates.
Its rejection of the same constants at $\lambda=1/20$ is a negative control
for the sufficient argument, not a counterexample to the theorem or to
the Berry--Esseen conjecture. The range $1/500<\lambda$ still needs a
different estimate or sharper constants.
